% ======================================================================
\subsection{PCP self-energy: irreducibility analysis}
\label{sub:pcp_irreducibility}
% ======================================================================

Starting from the verified-correct PCP self-energy (machine-precision
match to dense reference), we analyze every sum and product to determine
which can be collapsed, reordered, or shared, and which are irreducibly
expensive.  Every step below is a true algebraic identity.


% -------------------------------------------------------
\subsubsection{The exact expression}
% -------------------------------------------------------

\paragraph{Definitions.}
\begin{align}
  &\text{Shifted-FT mode:}\quad
  h_s^{\xi}(\vq, l)_{a}
  = \sum_{t=0}^{L-1} e^{-i\vq\cdot\vR_t}\,
    A_s^{\xi}\bigl((t{-}l)\bmod L,\; \kappa_a, \alpha_a\bigr)
  \in \CC^M,
  \label{eq:h_def}\\[6pt]
  &\text{External weight:}\quad
  w_a^{(s,\xi,l)}
  = A_s^{\xi}\bigl((-l)\bmod L,\; \kappa_a, \alpha_a\bigr)
  \in \RR,
  \label{eq:w_def}\\[6pt]
  &\text{Scalar projection:}\quad
  g_{s,s'}^{\xi\xi'}(\omega;\vq,l,l')
  = h_s^{\xi}(\vq,l)^\intercal\;
    G^{\gtrless}(\omega;\vq)\;
    h_{s'}^{\xi'\,*}(\vq,l'),
  \label{eq:g_def}
\end{align}
where $h^*$ denotes element-wise conjugation and $M = 3n_\mathrm{at}$
is the number of primitive-cell DOFs.

The notation $h^\intercal G h^*$ means $\sum_{c,f} h_c \, G_{cf} \, \overline{h_f}$.
This is the correct contraction from the self-energy
(the left vertex contributes $h_c$ and the conjugated right vertex
contributes $\overline{h_f}$).

\paragraph{Self-energy.}
\begin{equation}\label{eq:sigma_exact}
  \boxed{
  \Sigma^{\gtrless}_{ab}(\omega;\vq)
  = \frac{i\hbar}{2N_q}
    \sum_{\vq'}
    \sum_{\xi,\xi'}
    \frac{\lambda_\xi\lambda_{\xi'}}{36}
    \sum_{\sigma,\sigma'\in S_3}
    \sum_{l,l'=0}^{L-1}
    w_a^{(\sigma_1,\xi,l)}\;
    w_b^{(\sigma'_1,\xi',l')}\;
    C_{\sigma_1\sigma'_1}^{\xi\xi'}(\omega;\vq,\vq',l,l'),
  }
\end{equation}
where the convolution kernel is
\begin{equation}\label{eq:C_def}
  C_{\sigma_1\sigma'_1}^{\xi\xi'}(\omega;\vq,\vq',l,l')
  = \mathrm{IFFT}_\omega\!\bigl\{
      \hat{g}_{\sigma_2,\sigma'_3}^{\xi\xi'}(\vq',l,l')
      \;\hat{g}_{\sigma_3,\sigma'_2}^{\xi\xi'}(\vq{-}\vq',l,l')
    \bigr\},
\end{equation}
$\hat{g}$ denoting the FFT of~$g$ over~$\omega$.
Eqs.~\eqref{eq:sigma_exact}--\eqref{eq:C_def} have been
verified numerically: the shifted-FT PCP self-energy matches the
dense reference at $\sim 10^{-15}$ relative error on a $4\times 4$
$\vq$-mesh with a $2\times 2\times 2$ supercell (all~$\vq$ non-commensurate).


% -------------------------------------------------------
\subsubsection{Sum-by-sum analysis: what can be collapsed?}
\label{sssub:sums}
% -------------------------------------------------------

We examine each summation in \cref{eq:sigma_exact} and determine
whether it can be reduced.


\paragraph{1. The $\vq'$-sum (momentum convolution).}
$C$ depends on both $\vq'$ and $\vq{-}\vq'$, so the
$\vq'$-sum cannot be factored.
\textbf{However}, the accumulation into $\Sigma_{ab}$ is \emph{linear}
in $C$, and the weights $w_a, w_b$ do not depend on~$\vq'$.
Therefore:
\begin{equation}\label{eq:qprime_sum_first}
  \bar{C}^{\xi\xi'}_{\sigma_1\sigma'_1}(\omega;\vq,l,l')
  = \sum_{\vq'} C^{\xi\xi'}_{\sigma_1\sigma'_1}(\omega;\vq,\vq',l,l')
\end{equation}
can be computed \textbf{before} the outer-product accumulation.
This saves a factor of~$N_q$ in the accumulation phase.
\emph{This is the only reordering that reduces cost.}


\paragraph{2. The $S_3 \times S_3$ permutation sum (36 terms).}

For a permutation $\sigma = (\sigma_1, \sigma_2, \sigma_3) \in S_3$,
fixing $\sigma_1$ leaves exactly two choices for $(\sigma_2, \sigma_3)$:
the remaining pair in either order.  So the sum over $S_3$ decomposes
as $\sum_{\sigma_1=0}^{2} \sum_{(\sigma_2,\sigma_3)}$ with two inner
terms.  The $36$~products in the $S_3 \times S_3$ sum decompose into
$9$~bins indexed by $(\sigma_1, \sigma'_1)$, each receiving~$4$
products.

Can we reduce further?  Each bin's 4~products are:
\begin{align*}
  (\sigma_1{=}0):\; &
  \bigl\{(1,2),\;(2,1)\bigr\} \times
  \bigl\{(1,2),\;(2,1)\bigr\} \\[2pt]
  \text{Products:}\quad &
  g_{1,2'}\,g_{2,1'} + g_{1,1'}\,g_{2,2'}
  + g_{2,2'}\,g_{1,1'} + g_{2,1'}\,g_{1,2'} \\
  &= 2\bigl(g_{1,2'}\,g_{2,1'} + g_{1,1'}\,g_{2,2'}\bigr).
\end{align*}
(Primed indices denote $\sigma'$.)
So each bin has a \textbf{factor-of-2 symmetry}: swapping $(\sigma_2,\sigma_3)$
and simultaneously swapping $(\sigma'_2,\sigma'_3)$ gives the same product.
This reduces 36 products to~18 unique products.

\textbf{Can we reduce below 18?}
No.  The two distinct products
$g_{1,2'}g_{2,1'}$ and $g_{1,1'}g_{2,2'}$ involve different pairs
of $g$-projections, and these projections are generally different
complex numbers.  No further algebraic simplification is possible
without assuming special structure in~$G$.


\paragraph{3. The $(\xi, \xi')$-sum (PCP rank).}

The scalar projection $g^{\xi\xi'}$ depends on \emph{both}
$\xi$ (through $h_s^\xi$ on the left) and $\xi'$ (through
$h_{s'}^{\xi'}$ on the right).  In the convolution product:
\begin{equation}
  \hat{g}_{\sigma_2,\sigma'_3}^{\xi\xi'} \cdot
  \hat{g}_{\sigma_3,\sigma'_2}^{\xi\xi'},
\end{equation}
both factors share the \emph{same} $(\xi, \xi')$ pair.

\emph{Can the $(\xi,\xi')$ sum be separated?}
Write the product explicitly:
\[
  \bigl[h_{\sigma_2}^{\xi\,\intercal} \hat{G} h_{\sigma'_3}^{\xi'\,*}\bigr]
  \bigl[h_{\sigma_3}^{\xi\,\intercal} \hat{G} h_{\sigma'_2}^{\xi'\,*}\bigr].
\]
The first factor depends on $(\xi, \xi')$ through $(h^{\xi}, h^{\xi'})$,
and so does the second.  The product $f(\xi,\xi')\cdot g(\xi,\xi')$ is
\textbf{not} separable as $f(\xi)\cdot g(\xi')$ in general.

Therefore, the $(\xi,\xi')$-sum must be carried through the product,
giving $N_c^2$ terms.  \textbf{This is irreducible.}


\paragraph{4. The $(l, l')$-sum (cell shifts).}

This is the key overhead of the shifted-FT approach.
The question: \emph{does $g^{\xi\xi'}(\vq, l, l')$ have
exploitable structure in $(l, l')$?}

\subparagraph{Circular convolution structure.}
By \cref{eq:h_def}, $h_s^\xi(\vq, l)$ is related to $h_s^\xi(\vq, 0)$
by a circular shift of the real-space modes:
\begin{equation}\label{eq:h_shift}
  h_s^\xi(\vq, l)_a
  = \sum_t e^{-i\vq\cdot\vR_t}\, A_s^\xi\bigl((t{-}l)\bmod L\bigr)
  = \sum_{t'} e^{-i\vq\cdot\vR_{(t'+l)\bmod L}}\, A_s^\xi(t').
\end{equation}
Define the phase matrix $\Phi(\vq) \in \CC^{L\times L}$ with entries
$\Phi_{l,t'} = e^{-i\vq\cdot\vR_{(t'+l)\bmod L}}$.  Then
$h(\vq, l) = \Phi(\vq)_{l,:} \cdot \vec{A}$, i.e., $h$ is a
matrix-vector product where the row index is~$l$.

$\Phi$ is \textbf{not circulant} in general.
Its $(l, t')$-entry depends on $(t' + l) \bmod L$, but through
$\vR_{(t'+l)\bmod L}$, which is a lattice vector (not a linear
function of $t'+l$ when the modular reduction wraps around and
$\vq\cdot\vR$ is not periodic on the supercell lattice).

At \textbf{commensurate}~$\vq$ ($\vq\cdot\vR_L = 0 \bmod 2\pi$
for all lattice vectors $\vR_L$ of the supercell):
$\Phi_{l,t'} = e^{-i\vq\cdot\vR_{t'}} e^{-i\vq\cdot\vR_l}$,
so $\Phi$ becomes rank-1 (outer product of two phase vectors),
and $h(\vq, l) = e^{-i\vq\cdot\vR_l}\, h(\vq, 0)$.
The $(l,l')$-sum then collapses via $\sum_l e^{-i\vq\cdot\vR_l} w(l)
= \hat{f}(\vq)$.

At \textbf{non-commensurate}~$\vq$: $\Phi$ has full rank~$L$, so
$h(\vq, l)$ has $L$ linearly independent vectors as $l$ varies.
The projection $g(l, l') = h(l)^\intercal G\, h(l')^*$ is then a
general $L\times L$ matrix (not circulant, not low-rank).

\textbf{Conclusion:} \emph{At non-commensurate $\vq$, the $(l,l')$-sum
has $L^2$ truly independent terms.  The $L^2$ overhead is irreducible.}


\subparagraph{Fourier-$l$ transform does not help.}
One might try to diagonalize the $(l,l')$ structure via DFT:
\begin{equation}
  \hat{g}_{ss'}^{\xi\xi'}(\vq, k, k')
  = \sum_{l,l'} e^{-2\pi i(kl + k'l')/L}\;
    g_{ss'}^{\xi\xi'}(\vq, l, l').
\end{equation}
This factorizes as
$\hat{g}(k, k') = \hat{h}(\vq, k)^\intercal\, G\, \hat{h}(\vq, k')^*$,
where $\hat{h}(\vq, k) = \sum_l e^{-2\pi ikl/L} h(\vq, l)$.
The cost of computing $\hat{h}$ is $O(L \log L)$ per mode via FFT.
But the number of $(k, k')$~pairs is still~$L^2$, so the
projection cost is unchanged.  The FFT merely reorganizes where
the $L^2$~work is done (precomputation vs.\ projection), without
reducing the total.


% -------------------------------------------------------
\subsubsection{Optimal evaluation strategy}
\label{sssub:optimal}
% -------------------------------------------------------

Given the irreducibility results above, the optimal strategy is:

\paragraph{Phase 0 (once, before SCBA).}
Precompute $h_s^\xi(\vq, l) \in \CC^M$ for all
$(s, \xi, l, \vq)$: $3 N_c L N_q$ vectors of length~$M$.
Cost: $3 N_c M L \log L \cdot N_q$ (via circulant FFT).
Storage: $3 N_c L N_q M$ complex entries.

Also precompute external weights $w_a^{(s,\xi,l)}$ and
$\lambda_\xi \lambda_{\xi'} / 36$.

\paragraph{Phase 1 (per SCBA iteration): scalar projections.}
For each $(\vq, \omega)$: compute $G(\omega;\vq) \in \CC^{M \times M}$
(from the NEGF solver), then form the projections.

Optimal order for fixed $(\vq, \omega)$:
\begin{enumerate}
  \item For each $(l', s', \xi')$: compute
    $\tilde{G}_{l',s',\xi'} = G \cdot h_{s'}^{\xi'\,*}(\vq, l')
    \in \CC^M$.
    Cost: $3 N_c L \cdot M^2$.
  \item For each $(l, s, \xi)$ and each $(l', s', \xi')$:
    $g = h_s^{\xi}(\vq, l)^\intercal \cdot \tilde{G}$.
    Cost: $9 N_c^2 L^2 \cdot M$.
\end{enumerate}
Total per $(\vq, \omega)$:
\begin{equation}\label{eq:cost_proj_detail}
  \cW_\mathrm{proj}^{(\vq,\omega)}
  = 3 N_c M^2 L + 9 N_c^2 M L^2.
\end{equation}
Over all $(\vq, \omega)$: $N_q \cdot N_\mathrm{fft} \cdot \cW_\mathrm{proj}$.

\paragraph{Phase 2: FFT convolution.}
For each $(l, l')$ pair:
\begin{enumerate}
  \item FFT $g_{ss'}^{\xi\xi'}(\omega; \vq, l, l')$ over $\omega$
    for all $(\vq, s, s', \xi, \xi')$:
    $9 N_c^2 N_q$ FFTs of length $N_\mathrm{fft}$.
  \item For each $\vq_\mathrm{ext}$: gather $\hat{g}$ at $\vq'$ and
    $\vq{-}\vq'$ for all $N_q$ values of $\vq'$; compute the 18 unique
    products (from the $S_3\times S_3$ factor-of-2 symmetry); IFFT the
    9~bins; extract the physical frequency range; sum over~$\vq'$.
  \item Result:
    $\bar{C}^{\xi\xi'}_{\sigma_1\sigma'_1}(\omega; \vq, l, l')
    \in \CC^{N_c \times N_c \times 3 \times 3 \times N_\omega}$
    for each $\vq_\mathrm{ext}$.
\end{enumerate}
Cost per $(l, l')$:
\begin{equation}
  \cW_\mathrm{conv}^{(l,l')}
  = 9 N_c^2 N_q \cdot N_\mathrm{fft} \log N_\mathrm{fft}
  + N_q \cdot 18 N_c^2 N_q N_\mathrm{fft}
  + 9 N_c^2 N_q^2 \cdot N_\mathrm{fft} \log N_\mathrm{fft}.
\end{equation}
Over all $(l, l')$: multiply by~$L^2$.

\paragraph{Phase 3: accumulation.}
After summing $C$ over $\vq'$ (\cref{eq:qprime_sum_first}), for
each $(\vq, l, l')$ we accumulate:
\begin{equation}
  \Sigma_{ab}(\omega; \vq)
  \;\mathrel{+}=\;
  \sum_{s_1, s'_1 = 0}^{2}
  \sum_{\xi, \xi'}
  \frac{\lambda_\xi \lambda_{\xi'}}{36}\;
  w_a^{(s_1,\xi,l)}\;
  \bar{C}^{\xi\xi'}_{s_1 s'_1}(\omega; \vq, l, l')\;
  w_b^{(s'_1,\xi',l')}.
\end{equation}
This is 9 rank-1 updates of $M \times M$ matrices (one per $(s_1, s'_1)$
pair), with the $(\xi, \xi')$-sum absorbed into a matrix multiply:
$\Sigma \mathrel{+}= w_a^\intercal \cdot C_\mathrm{weighted} \cdot w_b$.

Cost per $(\vq, \omega, l, l')$: $9(N_c^2 N_\omega + N_c M N_\omega + M^2 N_\omega)$.
Over all: $L^2 N_q \cdot 9(N_c^2 + N_c M + M^2) N_\omega$.


% -------------------------------------------------------
\subsubsection{Total cost comparison}
\label{sssub:costs}
% -------------------------------------------------------

\paragraph{Notation.}
$M = 3 n_\mathrm{at}$, $L = N_\mathrm{cell}$,
$N_c$ = PCP rank, $N_q = n_\mathrm{kpts}$,
$N_\tau = 2(n_\mathrm{low} + N_\omega) \approx 2 N_\omega$.

\paragraph{Dense self-energy.}
\begin{equation}
  \cW_\mathrm{dense}
  = N_q^2 \bigl(M^4 N_\tau \log N_\tau + M^5 N_\omega\bigr).
\end{equation}
The first term is the FFT convolution of $G^{\gtrless}$ (element-wise
product of $M^2 \times M^2$ matrices in the frequency domain);
the second is the vertex contraction.

\paragraph{PCP shifted-FT.}
\begin{equation}\label{eq:cost_pcp_shifted}
  \cW_\mathrm{PCP}
  = \underbrace{N_q N_\tau (3 N_c M^2 L + 9 N_c^2 L^2 M)}_{\text{projections}}
  + \underbrace{L^2 \cdot 18 N_c^2 N_q^2 N_\tau}_{\text{products}}
  + \underbrace{L^2 \cdot 9 N_c^2 N_q^2 N_\tau \log N_\tau}_{\text{FFT/IFFT}}
  + \underbrace{L^2 N_q \cdot 9 (N_c M + M^2) N_\omega}_{\text{accumulation}}.
\end{equation}

\paragraph{Numerical comparison (Si, $M{=}6$, $L{=}8$, $N_q{=}16$,
$N_\omega{=}101$, $N_\tau{=}202$).}

\begin{center}
\renewcommand{\arraystretch}{1.3}
\begin{tabular}{l@{\quad}rrr}
  \toprule
  & Dense & PCP $N_c{=}4$ & PCP $N_c{=}24$ \\
  \midrule
  Projections
    & --- & $2.6 \times 10^8$ & $3.4 \times 10^9$ \\
  Products
    & --- & $9.6 \times 10^8$ & $3.5 \times 10^{10}$ \\
  FFT/IFFT
    & $4.4 \times 10^7$ & $7.1 \times 10^9$ & $2.6 \times 10^{11}$ \\
  Accumulation
    & $4.4 \times 10^6$ & $2.3 \times 10^8$ & $6.1 \times 10^9$ \\
  \midrule
  \textbf{Total}
    & $\mathbf{4.9 \times 10^7}$ & $\mathbf{8.3 \times 10^9}$ & $\mathbf{3.0 \times 10^{11}}$ \\
  \textbf{Ratio to dense}
    & $1.0\times$ & $\mathbf{170\times}$ & $6100\times$ \\
  \bottomrule
\end{tabular}
\end{center}

\paragraph{Key finding.}
At $N_c{=}4$, the shifted-FT PCP is $\sim\!170\times$ more expensive
than dense when counting FLOPs.  The earlier estimate of
$0.47\times$ was based on an erroneous cost model that neglected:
\begin{enumerate}
  \item The correct $L^2$~factor on the FFT/IFFT step (the
    convolution must be done $L^2$~times, not once).
  \item The $N_q$ factor in the product step (the gather-and-multiply
    over all $\vq'$ for each $\vq_\mathrm{ext}$ costs $N_q$ per
    iteration, giving $N_q^2 L^2$ total).
\end{enumerate}
The dense kernel benefits from its $M^4 N_\tau$ scaling with
$M{=}6$ being small.  The PCP overhead $L^2 N_c^2 = 1024$ for
$N_c{=}4$ far exceeds the dense $M^4 = 1296$ term \emph{when
multiplied by the $N_q^2$ factor in the convolution}.


% -------------------------------------------------------
\subsubsection{Why can't this be improved?}
\label{sssub:why_not}
% -------------------------------------------------------

Three fundamental obstacles prevent further simplification:

\begin{enumerate}
  \item \textbf{The $(l,l')$-sum is irreducible.}
    At non-commensurate~$\vq$, the shifted-FT modes $h(\vq, l)$
    span an $L$-dimensional space as $l$ varies.  The bilinear
    projection $g(l,l') = h(l)^\intercal G\, h(l')^*$ is a general
    $L \times L$ matrix.  No change of basis (including FFT over~$l$)
    reduces the number of independent entries below~$L^2$.

  \item \textbf{The $(\xi,\xi')$-coupling is irreducible.}
    The convolution product $g^{\xi\xi'}(\vq') \cdot g^{\xi\xi'}(\vq{-}\vq')$
    couples $\xi$ and $\xi'$ through the \emph{same} pair in both
    factors.  The $\xi$-sum and $\xi'$-sum cannot be separated:
    the product $f(\xi,\xi') g(\xi,\xi')$ is not of the form
    $h(\xi) k(\xi')$ in general.

  \item \textbf{The convolution is non-separable in $\vq'$.}
    The product $\hat{g}(\vq') \cdot \hat{g}(\vq{-}\vq')$ couples
    $\vq'$ and $\vq{-}\vq'$, preventing factorization of the
    $\vq'$-sum.
\end{enumerate}

\paragraph{Comparison with the SVD (separable) approach.}
The SVD decomposition $M_a = \sum_r F_r \otimes h_r$ produces
vectors $F_r \in \CC^{M \times 1}$ and $h_r \in \RR^{d_\mathrm{sc}}$
that are \emph{already} in the supercell-atom basis.  The gathering
matrix $T(\vq)$ is applied \emph{once} to each $h_r$, giving
$\tilde{h}_r(\vq) = T(\vq)\, h_r \in \CC^M$.  There is
\textbf{no cell-shift sum}---the supercell structure is absorbed
into the real-space vectors $h_r$.  The SVD cost is
$O(R^2 M N_q N_\tau)$, with no $L^2$ overhead.

The PCP decomposition, by contrast, factors the FC3 into
\emph{primitive-cell-indexed} modes $A_s^\xi(l, \kappa, \alpha)$,
requiring an explicit cell-shift sum to reconstruct the
supercell-indexed vectors.  This is the structural reason for
the $L^2$ overhead.

\paragraph{When is PCP still useful?}
If $N_c \ll R$ (i.e., the PCP rank is much lower than the SVD rank
needed for the same accuracy), the PCP decomposition may win
for \textbf{larger systems} where $M$ is large.  The PCP cost
scales as $L^2 N_c^2 M$, while the SVD costs $R^2 M^2$.  When
$L^2 N_c^2 < R^2 M$, i.e., $N_c < R\sqrt{M}/L$, PCP wins.
For Si with $M{=}6$, $L{=}8$, $R{=}24$: the break-even is
$N_c < 24 \cdot 2.4 / 8 \approx 7$.  At $N_c{=}4$ the PCP mode
count is lower, but the $L^2{=}64$ overhead in the convolution
step still makes it slower in absolute terms for this small system.

For larger systems ($M \gg 6$, $L$ moderate), the balance shifts
in favor of PCP because the $M^2$~factor in SVD projections
dominates.
\end{content>
